\section{Introduction}\label{sec:intro}

The suspension of the simulated F-34B is a spring, mass and damper system whose spring
constant and damping constant are generated from my student number, so the plant is
unknown before the experiment and cannot be looked up. The aim of this lab is to recover a
transfer function for that plant from measurements alone, which is what system
identification means in practice: choose inputs that make the unknown parameters visible in
the output, then read those parameters off the response.

Two experiments are used because each one exposes a different part of the model. A step
input shows the transient directly, so the steady state gain, the presence or absence of
oscillation, and the damped period can be read from a single log. A set of sinusoidal
inputs samples the frequency response, which fixes the corner frequency and the phase
behaviour, and gives an independent check on the values taken from the step. Where the two
disagree, that disagreement is itself a result and is reported as such.

Section~\ref{sec:modelling} derives the model structure from the equation of motion given
in the brief, and states which measurable feature of a response determines which parameter.
Section~\ref{sec:system_id} applies that to the measurements. Section~\ref{sec:validation}
simulates the identified model and compares it against the same data it was built from.
